\documentclass[]{article}
\usepackage{comment}

\begin{document}

We would like to thank the reviewers for their detailed feedback, which helps to strengthen our paper. Reviewer 1 and 3 recognize the work's "novel contribution to the field of LA research" and "important insights for building measurement models that support LA tools". This aligns closely with the goals of LA research and the aims of our paper. We believe that certain points of Reviewer 2 reflect a misunderstanding of our study's aim and methodology. In this rebuttal, we will clarify these points and unify them with the suggestions provided by Reviewer 3. 

We realize that the majority of the points raised by Reviewer 2 and 3 concern a lack of coherence between the different sections in the paper. We chose to focus on technical aspects of the measurement models that we investigated, by providing a mathematical explanation of the basic IRT model before introducing its extension (the IRTree model) which the paper focuses on. It is our impression that this extension is relatively novel to many learning platforms, necessitating a thorough explanation. Nonetheless, upon rereading the paper, we agree that this focus has not served coherence. This will be solved by separating the technical details from the background literature and explicitly reiterating the findings in conjunction with the conclusions. We think this simple adjustment will greatly improve the coherence by shifting the emphasis to the research questions, findings, and conclusions.

Reviewer 2 raises concerns about data selection. However, on the contrary, we argue that with such massive unstructured learning data with high ecological validity, it is precisely the careful selection that is crucial for generalizability to other systems and circumstances. Our data is used as an illustration of how this model can be used to estimate problem-skipping as a separate latent trait in ability measurement. This makes our findings more applicable for researchers and practitioners using other systems. 

The anonymization of many references in the reference list is also questioned. To clarify, the citations that are masked are not necessarily self-citations. They are masked, as recommended in the submission guidelines, because they involve the learning platform used in the study. While these studies are cited to demonstrate the breadth of research topics (all relevant for learning analytics) that have been conducted with the learning platform (section 3.3.1), the core references involved in the main line of reasoning are not self-citations (and therefore not masked). If the meta-reviewers are afraid that this section reflects an echoing of own ideas, this paragraph can be omitted without harming the central point of the research paper. 

We believe that the insights drawn from our work are valuable for understanding and optimizing learning within a real-world learning context. As this is our first submission to the LAK conference, we are especially thankful for the valuable feedback and look forward to engaging further with the LAK community. 



\begin{comment}

%% R1 ----
% The submission extends previous work on response scale- and achievement test data by comparing a ’unidimensional IRT model, which assumes that problem-skipping does not influence student ability differently from accuracy, to a multidimensional IRTree model, which explicitly accounts for differences in propensity to skip an item’, thus trying to contribute to the improvement of educational (online) assessment taking into account problem-skipping. The authors introduce the IRTree framework to estimate the latent correlation between two different response types in online learning data and show that problem-skipping should be measured separately from accuracy.
% The submission is appropriately situated concerning the field's current state, including sufficient coverage of relevant literature and theories. An introduction to previous research and gaps in the field justifies the need for the study, which makes a novel contribution to the field of LA research by trying to solve an important educational problem in the real-world context. The submission is written clearly, and the figures and tables contribute to a better understanding of the text.


%% R2 ----
% The paper aims to test  if problem-skipping and accuracy distinguishable processes in online learning and test the relationship between these two constructs. The paper is well structured, well-written and flows well. I have major concerns.
%First, the claims are strong and the methods are hardly clear, even more the dataset is carefully selected and therefore, one can’t generalize when selecting a dataset like that.
%Second, it is unclear how the authors came to the conclusions, or how they tested their hypothesis.
%Third, the introduction, and background are mostly technical with missing serious discussion of the problem or articulating the gap.
%Fourth, there is not a single mention of learning analytics in the article, only one reference, and I am not sure that this article even fits the conference of learning analytics.

%There are almost 40% of the citations masked, this limits the judgment of the argument and how it is supported. Also, the idea of that much self-citation does point either the author is echoing own ideas like in an echo chamber or these ideas have not been taken by anyone else. Have they been taken or endorsed by others, they should have been cited. It is a very odd situation anyway.
%There is a long introduction introducing Classical IRT models and IRTree models that make the paper technical and hard to read for the general audience. On the one hand, IRT models are rather common and you don’t use that space in discussing their mathematical basics and on the other hand, that space would have been better used in giving a proper review of the literature.
%There is no review of the existing literature, and almost a faint  mention of the gap of research, there is only a section on Item omission in online learning
%The data set as it is stated  “selected a specific subset of data from the Series game. First, we sampled data from the 2-year period between 01-09-2014 and 31-08-2016” is very old, and may not reflect current practices. But more importantly, the choice is based on the modeling reasons not for the sake of solving an educational problem or based on a research question. I am afraid that the results of the analysis are biased by these choices.
%I am not sure that the methods used can answer these research questions, or in fact the way the research questions are formulated are answerable.
%I have no idea how the authors came to the conclusion or empirically tested the main assumption of the paper, they don’t say, and the methods that they used are unclearly articulated and we are not told what led to this conclusion.
%Given that the data was selected with certain conditions it is unclear how this generalizes beyond these. The authors, for instance, chose only students with stable performance. This is of course, a special case, and later on, all the conclusions are generalized widely.

% R3 ----
%In this paper, the authors compared a unidimensional IRT model (which assumes that problem-skipping does not influence student ability differently from accuracy), to a multidimensional IRTree model, which accounts for differences in propensity to skip an item. By collecting data from a large-scale study using an online learning platform, they found evidence that problem-skipping and ability are distinguishable processes and previous measurement model that assumes that problem-skipping and accuracy are unrelated, might be overestimating students' ability, leading adaptive systems to assign overly difficult items.

%The research provides important insights for building measurement models that can support LA tools. I really enjoyed reading the result section. I have the following suggestions for the authors to improve the paper:

%1. Introduction section is relatively too short. Although, the authors motivates their work in Section 2.3, it would be good to reorganise some of the content to have the problem, motivation and their contributions listed clearly in the Introduction.
%2. Problem with anonymous references: The method section (mostly section 3.1) is quite hard to comprehend as majority of the references which have been mentioned to read for more details are anonymized. I can understand the double-blind process and respect author's judgement, however, I strongly feel some of those references could be easily reported in 3rd person (or can be included an anonymised copy as an electronic supplement document).
%3. Difficult to understand the IRTree model definitions: The Section 2.2 needs improvement. The equations are reported but their form is not well explained. Similar issues can be noted in equaltions 5,6,7, and 8. They are not well explained. (Minor: In equation 8, it should be alpha instead of "a").
%4. The results of the paper and the discussions seems not well aligned. The interesting results reported are not well discussed and therfore not backed with any related work (or previous literature), which makes the claims in the results less reliable. I might have missed in the discussion, but I couldn't find it cloasely linked to paper theme and results.

% Minor:
%the numbers reported in Section 4.1 and 4.2 are not fully supported with statistical tests.
%Overall, it's an interesting paper with good insights. I would urge the authors to rethink the flow of the paper for better readability and impact.

\end{comment}


\end{document}